% chapter14.tex -- Experimental Validation
% Based on the author's UGThesis.docx and experimental data from results.csv.

\chapter{Experimental Validation}
\label{ch:experiments}

In this chapter we present experimental results that validate the theoretical complexity bounds of our algorithm. 
We implemented the algorithm in C++ and tested it on 35 diverse biconnected plane graphs, ranging from simple cycles to complex graphs with up to 25 vertices and tens of millions of triangulations. 
The experiments confirm the \(O(1)\) amortised time per triangulation and \(O(n)\) space usage.

\section{Experimental Setup}
\label{sec:setup}

The algorithm was implemented in C++ using standard data structures:
\begin{itemize}
    \item The global chord set \(S\) was implemented as \texttt{std::unordered\_set} with a custom hash function for edges (pairs of vertex indices).
    \item Linked lists for \(L\), \(GS\), and \(VGS\) were implemented using \texttt{std::list} with pointers for constant -time removal.
    \item The two -pointer safe root finding algorithm was implemented as described in Chapter~\ref{ch:safe_root_algorithm}.
\end{itemize}

\textbf{Hardware:} AMD Ryzen 5 5600G (2 cores, 3.9 GHz, 8 GB RAM), Linux Mint 22.2, GCC 13.3.0 with \texttt{-O3} optimisation.

\textbf{Measurement:} Execution time was measured using \texttt{std::chrono::steady\_clock} with nanosecond precision. 
For each test case we ran the algorithm 5 times and report the median time. 
Memory usage was obtained from the Resident Set Size (RSS) via platform -specific queries (\texttt{/proc/self/smaps\_rollup} on Linux). 
Peak memory is reported as bytes per vertex.

\textbf{Test suite:} The 35 test cases include:
\begin{itemize}
    \item Simple cycles \(C_4\) to \(C_{15}\) (Catalan number growth).
    \item Grid graphs \(4\times4\) and \(5\times5\).
    \item Complex graphs with 11-18 vertices and multiple faces.
    \item Quadrilateral chains and other biconnected structures.
\end{itemize}
The full list with all measurements is provided in the supplementary material.

\section{Performance Results}
\label{sec:results}

Tables~\ref{tab:cycles} and~\ref{tab:complex} show representative results. 
The complete dataset (all 35 instances) is available online\footnote{\url{https://github.com/TAR2003/ThesisGraph}}.

\begin{table}[htbp]
    \centering
    \caption{Simple cycles: time per triangulation remains nearly constant.}
    \label{tab:cycles}
    \begin{tabular}{|l|r|r|r|r|}
        \hline
        \textbf{Cycle} & \textbf{Vertices} & \textbf{Triangulations} & \textbf{Time (s)} & \textbf{ns/triang} \\
        \hline
        \(C_4\) & 4 & 2 & 0.000003 & 1401 \\
        \(C_5\) & 5 & 10 & 0.000009 & 927 \\
        \(C_6\) & 6 & 68 & 0.000096 & 1406 \\
        \(C_7\) & 7 & 546 & 0.000414 & 758 \\
        \(C_8\) & 8 & 4,872 & 0.003040 & 624 \\
        \(C_9\) & 9 & 46,782 & 0.025897 & 554 \\
        \(C_{10}\) & 10 & 474,180 & 0.257996 & 544 \\
        \(C_{11}\) & 11 & 5,010,456 & 2.733858 & 546 \\
        \(C_{12}\) & 12 & 54,721,224 & 25.280202 & 462 \\
        \(C_{13}\) & 13 & 613,912,182 & 262.822566 & 428 \\
        \(C_{14}\) & 14 & 7,042,779,996 & 3064.941383 & 435 \\
        \(C_{15}\) & 15 & 82,329,308,040 & 26861.404167 & 326 \\
        \hline
    \end{tabular}
\end{table}

For cycles, the time per triangulation decreases slightly as \(n\) grows (from 1401 ns to 326 ns for \(C_{15}\)), which is consistent with the amortised constant behaviour: the constant factor improves because the overhead of building the root triangulation is spread over exponentially many triangulations. 
The total time scales linearly with the number of triangulations.

\begin{table}[htbp]
    \centering
    \caption{Complex biconnected graphs (larger instances).}
    \label{tab:complex}
    \begin{tabular}{|l|r|r|r|r|}
        \hline
        \textbf{Graph} & \textbf{Vertices} & \textbf{Triangulations} & \textbf{Time (s)} & \textbf{ns/triang} \\
        \hline
        grid \(4\times4\) & 16 & 4,613,504 & 1.884 & 408 \\
        grid \(5\times5\) & 25 & 81,920,000 & 46.257 & 565 \\
        complex & 13 & 1,282,136 & 0.596 & 465 \\
        graph V13F4 & 13 & 9,568,244 & 4.438 & 464 \\
        graph V13F4 v2 & 13 & 10,225,274 & 4.537 & 444 \\
        graph V15F5 & 15 & 89,897,864 & 40.055 & 446 \\
        graph V18F6 & 18 & 7,932,840,664 & 3637.363 & 459 \\
        \hline
    \end{tabular}
\end{table}

The per -triangulation times for complex graphs (408-565 ns) are similar to those for cycles, confirming that the overhead of the global chord set and pruning does not degrade performance.

\section{Time Complexity Validation}
\label{sec:time_validation}

Figure~\ref{fig:loglog} plots the total running time against the number of triangulations on a log -log scale. 
The data points lie close to a straight line with slope 1, indicating that time is proportional to the number of triangulations. 
This empirical evidence strongly supports the theoretical claim of \(O(1)\) amortised time per triangulation.

\begin{figure}[htbp]
    \centering
    \includegraphics[width=0.8\textwidth]{total_time_vs_triangulations_median.png}
    \caption{Log -log plot of total running time (median) versus number of triangulations. 
    The linear fits slope=1 confirms that the algorithm runs in time proportional to the output size.}
    \label{fig:loglog}
\end{figure}

\section{Space Complexity Validation}
\label{sec:space_validation}

Table~\ref{tab:memory} shows peak memory usage for selected graphs. 
Memory per vertex ranges from 315 B to 1 KB, independent of the number of triangulations. 
For instance, the 13 -vertex complex graph (1.28 million triangulations) uses only 315 B/vertex, while the 25 -vertex grid (82 million triangulations) uses 655 B/vertex. 
This confirms the linear space bound \(O(n)\).

\begin{table}[htbp]
    \centering
    \caption{Memory usage (peak).}
    \label{tab:memory}
    \begin{tabular}{|l|r|r|r|}
        \hline
        \textbf{Graph} & \textbf{Vertices} & \textbf{Triangulations} & \textbf{Memory (B/vertex)} \\
        \hline
        \(C_5\) & 5 & 10 & 819 \\
        \(C_6\) & 6 & 68 & 683 \\
        complex & 13 & 1,282,136 & 315 \\
        graph V11F3 & 11 & 672,902 & 372 \\
        graph V12F3 (some) & 12 & 10,451,550 & 341 \\
        graph V15F5 & 15 & 89,897,864 & 546 \\
        grid \(5\times5\) & 25 & 81,920,000 & 655 \\
        \hline
    \end{tabular}
\end{table}

The slightly higher memory for very small graphs (e.g., \(C_5\): 819 B/vertex) is due to fixed overhead of data structures; for larger graphs the overhead becomes negligible.

\section{Discussion}
\label{sec:discussion}

The experimental results are fully consistent with the theoretical analysis:

\begin{itemize}
    \item \textbf{Constant amortised time:} The per -triangulation time varies between 300 and 1500 ns across all instances, with no systematic increase as the output size grows from 2 to over 80 million. 
    The small variations are due to differences in the constant factors (e.g., cycle length, number of faces, and the cost of hash set operations).
    
    \item \textbf{Linear space:} Memory usage remains between 300 B and 1 KB per vertex, even when the output contains tens of millions of triangulations. 
    This demonstrates that the algorithm does not store the entire list of triangulations; it generates them on the fly and discards them after output.
    
    \item \textbf{Handling separation pairs:} The test graphs include several biconnected graphs with separation pairs (e.g., the quadrilateral chains and the “complex” graph). 
    The algorithm successfully generates all triangulations without producing multi -edges, validating the pruning strategy and the safe root selection.
\end{itemize}

\section{Limitations}
\label{sec:limitations_experiment}

The current implementation uses C++ standard hash sets, which have expected constant time but may occasionally incur rehashing overhead. 
For graphs with more than 100 vertices, the number of triangulations becomes astronomically large, so we did not test beyond 25 vertices. 
Nevertheless, the asymptotic behaviour is clear from the growth pattern.

\section{Summary}
\label{sec:experiment_summary}

We have implemented the algorithm and tested it on 35 biconnected plane graphs. 
The experimental results confirm the theoretical bounds: \(O(1)\) amortised time per triangulation and \(O(n)\) space. 
The algorithm scales gracefully from small cycles (2 triangulations) to large grids (over 80 million triangulations), demonstrating its practical efficiency.